% !TEX root = ../notes_dgxspark.tex
% Chapter 3: Run on the GPU -- without taking all of it
%%%%%%%%%%%%%%%%%%%%%%%%%%%%%%%%%%%%%%%%%%%%%%%%%%%%%%%%%%%%%%%%%%%%%%
%
% The Spark has no dedicated GPU memory: CPU and GPU share 128 GB, and so
% do all users. nvidia-smi therefore prints "Memory-Usage: Not Supported"
% (NVIDIA known issue) but still lists per-process GPU memory; total use
% comes from free -h. The default cap of a quarter is Mark's call.

\index{GPU}
\chapter{Run on the GPU}
\printchapterversion{03_gpu}

\section{One Memory for Everyone}

\index{unified memory}\newthought{CPU and GPU share the same 128~GB.} Every
gigabyte your training takes is missing for everyone else, whether it sits
on the CPU side or the GPU side.
\begin{verbatim}
nvidia-smi    # who runs on the GPU, per process
free -h       # how much memory is left overall
\end{verbatim}
\marginnote{\newthought{Memory-Usage: Not Supported} in \texttt{nvidia-smi}
is expected. There is no separate GPU memory to report. The per-process
list below it still works.}

\section{Cap Your Framework}

\index{memory cap}\newthought{Frameworks take more than they need unless
told otherwise.} Cap yours at a quarter of the memory, 32~GB, before the
first tensor touches the GPU:
\begin{verbatim}
torch.cuda.set_per_process_memory_fraction(0.25)  # PyTorch
export XLA_PYTHON_CLIENT_MEM_FRACTION=0.25        # JAX
\end{verbatim}
\marginnote{\newthought{JAX grabs 75\%} of the memory at startup by
default, whether it needs it or not.}
\marginnote{\newthought{Need more than a quarter?} Ask me first.}

\section{Start Small, Then Scale}

\index{batch size}\index{bfloat16}\index{checkpoint}\newthought{Run a few
steps on a slice of the data first.} Watch \mbox{\texttt{free -h}} while it runs,
then raise the batch size until you reach your cap, not the machine's.
Train in \texttt{bfloat16}, reach a large effective batch through gradient
accumulation, and save checkpoints to \texttt{\textasciitilde/runs} so a
stopped job resumes instead of starting over.
\marginnote{\newthought{bfloat16 in PyTorch.} Wrap the training step in
\texttt{torch.autocast} with \texttt{dtype=torch.bfloat16}.}

\newthought{Memory is freed only when your process ends.} An idle Python
session or notebook kernel keeps holding it. Find yours in
\texttt{nvidia-smi} and end it.
